\section{Requerimientos Específicos}

La priorización de los componentes del sistema se ha abordado desde dos dimensiones: una estratégica basada en la técnica \textbf{MoSCoW} aplicada a las Historias de Usuario (ver Sección 5), y una técnica de impacto \textbf{Alto, Medio y Bajo} aplicada a los requisitos atómicos detallados en el \textbf{Apéndice A}. Esta sección se centra en las restricciones de diseño, interfaces y atributos de calidad que garantizan la viabilidad técnica de la solución.

\subsection{Interfaces Externas}

\subsubsection{Interfaces de Usuario}

\begin{itemize}
    \item \textbf{Portal Web:} Interfaz administrativa y de gestión basada en estándares \textbf{HTML5/CSS3}, optimizada para resoluciones Full HD (1920x1080).
    \item \textbf{Aplicación Móvil:} Interfaz táctil para dispositivos \textbf{Android 10+} con soporte nativo para gestos (pinch-to-zoom) y operación en entornos de baja luminosidad (\textbf{HPHMI Grayscale}).
\end{itemize}
\subsubsection{Interfaces de Hardware}

\begin{itemize}
    \item \textbf{Captura de Datos:} Integración con la cámara del dispositivo móvil para escaneo de códigos QR y captura de evidencia fotográfica.
    \item \textbf{Telemetría IoT:} Interfaz de recepción de datos para sensores de vibración, temperatura y presión mediante protocolos de baja latencia.
\end{itemize}
\subsubsection{Interfaces de Software}

\begin{itemize}
    \item \textbf{Motor de Visualización:} Soporte nativo para renderizado de mallas 3D en formatos \textbf{GLB/GLTF} y \textbf{STEP} simplificado.
    \item \textbf{Servicios IA:} Interfaz de comunicación con modelos de lenguaje (LLM) de clase industrial para el procesamiento de lenguaje natural (RAG).
\end{itemize}
\subsubsection{Interfaces de Comunicación}

\begin{itemize}
    \item \textbf{Seguridad de Red:} Todas las transacciones deben realizarse sobre \textbf{HTTPS} con cifrado \textbf{TLS 1.3}.
    \item \textbf{Protocolo Industrial:} Uso de \textbf{MQTT} para la transmisión de datos de sensores, optimizando el consumo de batería y ancho de banda en campo.
\end{itemize}
\subsection{Requisitos No Funcionales (Atributos de Calidad)}

Basados en el estándar \textbf{ISO/IEC 25010}, se definen los siguientes objetivos de calidad:

\begin{itemize}
    \item \textbf{Seguridad (Security):} Implementación de cifrado \textbf{AES-256} en reposo para documentos técnicos y cadena de integridad criptográfica para la auditoría.
    \item \textbf{Eficiencia de Desempeño (Performance):} Tiempo de respuesta de la visualización 3D $\ge$ 30 FPS en dispositivos móviles y latencia de consulta IA $<$ 15 segundos.
    \item \textbf{Fiabilidad (Reliability):} Capacidad de operación \textbf{Offline-First} con sincronización automática y manejo de conflictos de datos.
    \item \textbf{Mantenibilidad (Maintainability):} Arquitectura basada en \textbf{Requisitos Transversales (TR)} que permite la actualización de patrones de seguridad sin afectar la lógica de negocio.
\end{itemize}
\subsection{Patrones de Requisitos Transversales (TR)}

El sistema implementa 11 patrones arquitectónicos comunes que garantizan la integridad y escalabilidad de la solución:

\begin{enumerate}
    \item \textbf{TR-001 (Auditoría Inmutable):} Registro inalterable de toda acción crítica mediante tecnología \textit{Append-only}.
    \item \textbf{TR-002 (Cifrado en Reposo):} Protección de datos almacenados mediante el estándar \textbf{AES-256}.
    \item \textbf{TR-003 (Integridad de Datos):} Verificación de la autenticidad de archivos y registros mediante \textbf{Hashing SHA-256}.
    \item \textbf{TR-004 (Limitación de Tasa):} Control de tráfico (\textit{Rate Limiting}) para prevenir abusos en APIs y servicios de IA.
    \item \textbf{TR-005 (Cifrado en Tránsito):} Aseguramiento de las comunicaciones mediante protocolos \textbf{TLS 1.3}.
    \item \textbf{TR-006 (Gobernanza RBAC):} Control de acceso basado en roles industriales y principio de menor privilegio.
    \item \textbf{TR-007 (Sincronización Offline):} Gestión de datos y persistencia local para zonas sin conectividad de red.
    \item \textbf{TR-008 (ISO 14224):} Gestión y validación de jerarquía de activos industriales (9 niveles).
    \item \textbf{TR-009 (Framework de IA/LLM):} Gestión unificada de prompts, contextos RAG y \textit{scoring} de confianza.
    \item \textbf{TR-010 (Sincronización en Tiempo Real):} Sincronización instantánea de capas visuales mediante \textbf{WebSockets}.
    \item \textbf{TR-011 (Cumplimiento HPHMI):} Estándar visual de alto desempeño para reducción de carga cognitiva.
\end{enumerate}
